\section{A sharp integral lift in the Wach lattice at each cut}

% BEGIN R82W1
\paragraph{R82W1. The least power on an individual letter.}
On an unramified $K$ at odd $p$ and level one use the integral Wach
lattices of W79H1--W79H5. For an onto map $f:T_1\to T_2$ of
crystalline lattices, suppose the labelled jumps of the source on
letter $i$ lie in $[-M_i,h_i]$, with $M_i,h_i\ge0$. Put
$H_i=M_i+h_i$, $H_* =\max_i H_i$, $A=O[[X]]$,
$Q=((1+X)^p-1)/X$. Write $N_i=N_i(T_2)$, $N'_i=fN_i(T_1)$
and $I_i=\mathrm{Ann}_A(N_i/N'_i)$. All letters are indexed in
Frobenius order. There is a least $s_i\ge0$ with $X^{s_i}\in I_i$,
and $s_i\le H_*$.

To check this with varying heights, strictness gives the same
interval on the target. Modulo every coefficient-uniformizer
power, A81R8 gives a common $X^{b_d}$ killing the cokernels
after clearing the localized surjections. Its coefficient
extraction uses $Q^{M_i}\varphi N_{i-1}\subset N_i$ on the target
and $Q^{h_i}N'_i\subset A\varphi N'_{i-1}$ on the source image.
Both follow from W79H4 at that letter. In the expansion
$Q^{H_i}a=\sum_{k<p}(1+X)^k\varphi(a_k)$ for any annihilator
$a$ at $i$, uniqueness in the Frobenius basis therefore gives
$a_k$ in the annihilator at $i-1$. This also holds for
$Q^{H_*}a$ by first multiplying $a$ by $Q^{H_*-H_i}$.
Thus the digit extractions of A81R8 on $X^{H_*}(\epsilon-1)^{b_d}$
with $\epsilon=1+X$ give $X^{H_*}$ in the common annihilator modulo every depth.
Krull intersection on each finite cokernel gives $X^{H_*}\in I_i$.
The same coefficient test holds integrally with $H_i$: apply it at
every depth and use Krull intersection. In particular
\[
       b\ge H_i,\quad X^bQ^{b-H_i}\in I_i
                      \quad\Longrightarrow\quad X^b\in I_{i-1},
\]
since $X^bQ^b=\varphi(X^b)$ has just its constant Frobenius-basis
component. This implication retains the local height $H_i$.
% END R82W1

% BEGIN R82W2
\paragraph{R82W2. Descending one power of the series parameter.}
Let $q=p-1$, and take $\gamma$ with cyclotomic value
$u=\omega(1+p)$ where $\omega$ generates the Teichm\"uller
units of order $q$. This element exists in the Galois quotient
over the unramified base. On both $N_i$ and $N'_i$ the action
is the identity modulo $X$ in the sense
$(\gamma-1)N_i\subset XN_i$ and $(\gamma-1)N'_i\subset XN'_i$.
Indeed in the regular Taylor lattice at zero all constant
crystalline periods are fixed under this action; evaluation at zero
and integral divisibility prove the first assertion. Apply the same
fact to the source and map forward for the second.
Also $\gamma(X)/X=u+Xw(X)$ with integral $w$.

If $k>0$ and $p^aX^kN_i\subset N'_i$, apply $\gamma-1$ to
$p^aX^kz$ for an arbitrary $z\in N_i$ and divide by $X$.
The result lies in $N'_i$ and is
$p^a(u^k-1)X^{k-1}z$ plus an element of $p^aX^kN_i\subset N'_i$.
Consequently $p^{a+v_p(u^k-1)}X^{k-1}\in I_i$.
Starting with $X^{s_i}$ gives
\[
  p^{D(j,s_i)}X^j\in I_i,\qquad
  D(j,s)=\sum_{j<k\le s}v_p(u^k-1)\quad(0\le j\le s).
\]
Only multiples of $q$ contribute. There the valuation is
$1+v_p(k)$ by the binomial formula. Thus $D$ counts multiples
of $qp^v$ in $(j,s]$, summed over $v\ge0$. There are at most
$1+\log_p(H_*+1)$ nonempty summands. Bounding each count by
$(s-j)/(qp^v)+1$ gives
\[
 D(j,s)\le (s-j)p/q^2+c_H,\qquad
 c_H=\lceil 1+\log_p(H_*+1)\rceil .
\]
For $j\ge s_i$ the $X$-power alone suffices. In particular this
argument bounds every vector of the target before taking tame
invariants, and its cost is uniform in the ranks.
% END R82W2

% BEGIN R82W3
\paragraph{R82W3. The forward weighted height sum.}
Set $s=s_i$ and $b=H_i+c_H+\lceil s/q\rceil$. We prove
$X^bQ^{b-H_i}\in I_i$. Write $Q=X^q+pV$ with $V\in A$.
Each binomial summand is an integral multiple of
$p^aX^{k(a)}$, where $0\le a\le b-H_i$ and
$k(a)=b+q(b-H_i-a)$. For $k(a)\ge s$ it is already in $I_i$.
Otherwise it is enough by R82W2 to have
$a\ge(s-k(a))p/q^2+c_H$. The right side minus $a$
is increasing in $a$, with slope $p/q-1>0$. At the upper
endpoint $a=b-H_i$ we have $k(a)=b$. Now
$b-H_i-c_H\ge s/q$ and $b\ge s/q$, so
\[
 (s-b)p/q^2\le s(q-1)p/q^3\le s/q
\]
using $(q-1)p=q^2-1$. This proves the required bound for all
summands. R82W1 therefore yields
$s_{i-1}\le H_i+c_H+1+s_i/q$. Iterating and using
$s_k\le H_*$ on the cyclic word makes the final geometric
remainder tend to zero. In particular
\[
   s_i\le\min\left\{H_*,
    \left\lfloor\sum_{a\ge0}
         \frac{H_{i+a+1}+c_H+3}{q^a}\right\rfloor\right\}.
\]
The two extra units in the numerator put the statement exactly
in the convention of $\mathrm{sBd}$. The shift $i+a+1$ is
essential when a cut is followed by a string of thin letters.
% END R82W3

% BEGIN R82W4
\paragraph{R82W4. The two cut lifts at the prescribed exponent.}
Take $j=0$ in R82W2 and insert R82W3's upper bound $B_i$ for
$s_i$. Since the cost is integral, $p^{r_F(i)}N_i\subset N'_i$
with $r_F(i)=\lfloor B_i p/q^2\rfloor+c_H$, exactly the
source's $rF$ for this height sequence. If the target vector is
fixed under the tame group $\Delta$, take any preimage and average
over $\Delta$. Its order $p-1$ is a unit and $f$ commutes with the
action, giving an invariant preimage of that same multiple.
This proves $\mathrm{WPrimeCuts}$ on the invariant coordinates at
all cuts, since the uniform $rr_{\rm par}$ is the maximum over
the cuts and over the two tested height sequences of these $r_F$.

Apply the same assertion to the integral extension representing
an arbitrary finite target class, projected onto its trivial
quotient. Its rational extension is crystalline by the definition
of finiteness of the class. The additional zero jump lies in each
$[-M_i,h_i]$ of the target, so the other of the two $r_F$ sequences
gives at each cut an invariant Wach lift of $p^{rr_{\rm par}}$
times the quotient generator. This proves both the map-lifting
and the quotient-generator assertions at the actual cut exponent
for I82W4's word. These assertions concern vectors in the Wach
lattice. Lifting a finite cohomology class also requires compatible
Frobenius and generator differences along the arcs.
% END R82W4
